% c02 原书 p.10–18（PDF p.017–025）
\begin{proposition}\label{it:1.1.4}
设 $X$ 是拓扑空间。则 $X$ 上交换群的前层范畴是阿贝尔范畴。设 $\phi\colon\cF\to\cG$ 是交换群前层的态射。则 $\phi$ 的核、余核与像是如下定义的前层：
\[
\begin{aligned}
(\ker\phi)(U)&=\ker\bigl(\phi(U)\colon\cF(U)\to\cG(U)\bigr),\\
(\coker\phi)(U)&=\coker\bigl(\phi(U)\colon\cF(U)\to\cG(U)\bigr),\\
(\operatorname{im}\phi)(U)&=\operatorname{im}\bigl(\phi(U)\colon\cF(U)\to\cG(U)\bigr),
\end{aligned}
\]
其中 $U$ 取遍 $X$ 的开子集。这些前层在点 $P\in X$ 处的茎由
\[
\begin{aligned}
(\ker\phi)_P&\cong\ker(\phi_P\colon\cF_P\to\cG_P),\\
(\coker\phi)_P&\cong\coker(\phi_P\colon\cF_P\to\cG_P),\\
(\operatorname{im}\phi)_P&\cong\operatorname{im}(\phi_P\colon\cF_P\to\cG_P)
\end{aligned}
\]
给出。
\end{proposition}

\begin{proposition}\label{it:1.1.5}
设 $\cF$ 是拓扑空间 $X$ 上的前层。则存在一个由层 $\cF^+$ 与态射 $\theta\colon\cF\to\cF^+$ 组成的对 $(\cF^+,\theta)$，使得对任意层 $\cG$ 与任意态射 $\phi\colon\cF\to\cG$，存在唯一的态射 $\psi\colon\cF^+\to\cG$，使 $\phi=\psi\theta$。对 $(\cF^+,\theta)$ 在相差唯一同构的意义下唯一。对任意点 $P\in X$，$\theta_P\colon\cF_P\to\cF^+_P$ 是同构。我们称 $\cF^+$ 为前层 $\cF$ 所伴随的层。
\end{proposition}

\begin{proof}
对 $X$ 的任意开子集 $U$，定义 $\cF^+(U)$ 为满足以下两个条件的函数 $s\colon U\to\coprod_{P\in X}\cF_P$ 构成的集合：
\begin{enumerate}[label=(\alph*),leftmargin=2.8em]
  \item 对每个 $P\in U$，有 $s(P)\in\cF_P$。
  \item 对每个 $P\in U$，存在 $P$ 的含于 $U$ 的邻域 $U_P$ 与截面 $t\in\cF(U_P)$，使得对任意 $Q\in U_P$，$s(Q)$ 是 $t$ 在 $Q$ 处的芽 $t_Q$。
\end{enumerate}
则 $\cF^+$ 是层，且有典范态射 $\theta\colon\cF\to\cF^+$。该态射对每个 $P\in X$ 诱导茎上的同构 $\theta_P\colon\cF_P\to\cF^+_P$。对每个 $s\in\cF^+(U)$，存在 $U$ 的开覆盖 $\{U_i\}_{i\in I}$ 与截面 $t_i\in\cF(U_i)$，使 $s|_{U_i}=\theta(t_i)$。对每个 $P\in U_i\cap U_j$，有
\[
t_{i,P}=s(P)=t_{j,P}.
\]
取 $P$ 在 $U_i\cap U_j$ 中的开邻域 $U_P$，使 $t_i|_{U_P}=t_j|_{U_P}$。则这些 $U_P$（$P\in U_i\cap U_j$）构成 $U_i\cap U_j$ 的开覆盖，且 $\phi(t_i)|_{U_P}=\phi(t_j)|_{U_P}$。因 $\cG$ 是层，必有 $\phi(t_i)|_{U_i\cap U_j}=\phi(t_j)|_{U_i\cap U_j}$。因 $\{U_i\}_{i\in I}$ 构成 $U$ 的开覆盖且 $\cG$ 是层，可找到截面 $\psi(s)\in\cG(U)$，使 $\psi(s)|_{U_i}=\phi(t_i)$ 对所有 $i$ 成立。截面 $\psi(s)$ 由以下性质唯一刻画：
\[
\psi(s)_P=\phi_P(s(P))
\]
对所有 $P\in U$ 成立，其中 $\phi_P\colon\cF_P\to\cG_P$ 是 $\phi$ 诱导的茎上的态射。可以验证 $\psi$ 定义了一个层态射，且具有所需的性质。
\end{proof}
\begin{tnote}
\emph{层化}\idx{cenghua}{层化}的构造中，条件 (b) 不可省。只要求 $s(P)\in\cF_P$（条件 (a)），得到的只是任意选取的一族芽，它一般不是层：既不能保证 $s$ 局部来自同一个截面，也不能保证 $\theta_P$ 是同构。条件 (b) 正是要求 $s$ 局部来自 $\cF$ 的截面，这是 $\cF^+$ 成为层、且 $\theta_P$ 为同构的关键。直观上，$\cF^+$ 就是把 $\cF$ 的截面按「局部相等」等同之后得到的最小层。[依据：原书 p.10 命题 1.1.5 的证明与 p.11 第 1 段]
\end{tnote}

注意，对每个截面 $s\in\cF^+(U)$，都可找到 $U$ 的开覆盖 $\{U_i\}_{i\in I}$ 与截面 $s_i\in\cF(U_i)$，使 $\theta(s_i)=s|_{U_i}$；也就是说，$\cF^+$ 中的截面局部都来自 $\cF$ 的截面。两个截面 $s,t\in\cF(U)$ 在 $\cF^+(U)$ 中有相同的像，当且仅当存在 $U$ 的开覆盖 $\{U_i\}_{i\in I}$，使 $s|_{U_i}=t|_{U_i}$ 对所有 $i$ 成立。于是 $\cF$ 中局部相等的两个截面在 $\cF^+$ 中被等同。当 $\cF$ 是层时，$\theta\colon\cF\to\cF^+$ 是同构。

\begin{proposition}\label{it:1.1.6}
设 $X$ 是拓扑空间。则 $X$ 上交换群的层范畴是阿贝尔范畴。设 $\phi\colon\cF\to\cG$ 是交换群层的态射。则 $\ker\phi$ 是由
\[
(\ker\phi)(U)=\ker(\phi(U))
\]
定义的层，$\coker\phi$ 是前层 $U\mapsto\coker(\phi(U))$ 所伴随的层，而 $\operatorname{im}\phi$ 是前层 $U\mapsto\operatorname{im}(\phi(U))$ 所伴随的层。
此外，对任意 $P\in X$，有
\[
\begin{aligned}
(\ker\phi)_P&\cong\ker(\phi_P),\\
(\coker\phi)_P&\cong\coker(\phi_P),\\
(\operatorname{im}\phi)_P&\cong\operatorname{im}(\phi_P).
\end{aligned}
\]
\end{proposition}

\begin{proof}
可以证明，由 $U\mapsto\ker(\phi(U))$ 定义的前层是层，且是 $\phi$ 的核。利用 \ref{it:1.1.4} 与 \ref{it:1.1.5}，可以证明 $\phi$ 的余核是前层 $U\mapsto\coker(\phi(U))$ 所伴随的层。记 $\cC$ 为由 $U\mapsto\coker(\phi(U))$ 定义的前层，记 $\cP$ 为由 $U\mapsto\operatorname{im}(\phi(U))$ 定义的前层。我们来证明 $\cP^+$ 是 $\phi$ 的像，即 $\phi$ 的余核的核。有前层的典范态射 $\cP\to\cG$，它诱导态射 $\cP^+\to\cG$。因复合 $\cP\to\cG\to\cC$ 为 $0$，复合 $\cP\to\cG\to\cC^+$ 也为 $0$，从而复合 $\cP^+\to\cG\to\cC^+$ 为 $0$。于是 $\cP^+\to\cG$ 诱导态射 $\cP^+\to\ker(\cG\to\cC^+)$。我们断言它是同构。由 \ref{it:1.1.1}，只需证明对每个 $P\in X$，同态 $(\cP^+)_P\to(\ker(\cG\to\cC^+))_P$ 是同构。由 \ref{it:1.1.4} 与 \ref{it:1.1.5}，有
\begin{align*}
(\cP^+)_P&\cong\cP_P\cong\operatorname{im}(\phi_P),\\
(\ker(\cG\to\cC^+))_P&\cong\ker\bigl(\cG_P\to(\cC^+)_P\bigr)\cong\ker(\cG_P\to\cC_P)\\
&\cong\ker\bigl(\cG_P\to\coker(\phi_P)\bigr).
\end{align*}
断言随即成立。命题中其余断言的证明留给读者。
\end{proof}

\begin{lemma}\label{it:1.1.7}
设 $f\colon A\to B$ 是阿贝尔范畴中的态射。则
\begin{enumerate}[label=(\roman*),leftmargin=2.8em]
  \item $f$ 是单射，当且仅当 $\ker f=0$；
  \item $f$ 是满射，当且仅当 $\coker f=0$，当且仅当 $\operatorname{im}f=B$。
\end{enumerate}
\end{lemma}

\begin{proof}
我们只证明 (i) 中「当」的部分，其余留给读者。考虑交换图
\begin{center}
\begin{tikzcd}
A \arrow[r] \arrow[d] & B \\
\operatorname{coim}f \arrow[r] & \operatorname{im}f \arrow[u]
\end{tikzcd}
\end{center}
若 $\ker f=0$，则左边竖直箭头是同构。下边水平箭头是同构，而右边竖直箭头是单态射。所以上边水平箭头是单态射。
\end{proof}

\begin{corollary}\label{it:1.1.8}
设 $X$ 是拓扑空间。
\begin{enumerate}[label=(\roman*),leftmargin=2.8em]
  \item 在 $X$ 上交换群的前层范畴中，态射 $\phi\colon\cF\to\cG$ 是单射（resp.\ 满射，resp.\ 双射），当且仅当对 $X$ 的每个子集 $U$，$\phi(U)\colon\cF(U)\to\cG(U)$ 是单射（resp.\ 满射，resp.\ 双射）。
  \item 在 $X$ 上交换群的层范畴中，态射 $\phi\colon\cF\to\cG$ 是单射，当且仅当对 $X$ 的每个开子集 $U$，$\phi(U)\colon\cF(U)\to\cG(U)$ 是单射。
  \item 在 $X$ 上交换群的层范畴中，态射 $\phi\colon\cF\to\cG$ 是单射（resp.\ 满射，resp.\ 双射），当且仅当对每个 $P\in X$，$\phi_P\colon\cF_P\to\cG_P$ 是单射（resp.\ 满射，resp.\ 双射）。
  \item 在 $X$ 上交换群的层范畴中，态射 $\phi\colon\cF\to\cG$ 是满射，当且仅当对任意开子集 $U$ 与任意截面 $t\in\cG(U)$，存在 $U$ 的开覆盖 $\{U_i\}_{i\in I}$ 与截面 $s_i\in\cF(U_i)$，使 $\phi(s_i)=t|_{U_i}$ 对所有 $i$ 成立。
\end{enumerate}
\end{corollary}

\begin{proof}
(i) 由 \ref{it:1.1.4} 与 \ref{it:1.1.7} 得到。(ii) 由 \ref{it:1.1.6} 与 \ref{it:1.1.7} 得到。(iv) 由 (iii) 得到。下面证明 (iii) 中关于满射的断言，其余留给读者。由 \ref{it:1.1.7}，$\phi$ 是满射当且仅当 $\operatorname{im}\phi\cong\cG$。由 \ref{it:1.1.1}，这等价于说 $(\operatorname{im}\phi)_P\cong\cG_P$ 对每个 $P\in X$ 成立。由 \ref{it:1.1.6}，这等价于说 $\operatorname{im}(\phi_P)\cong\cG_P$，即 $\phi_P$ 对每个 $P\in X$ 是满射。
\end{proof}
\begin{tnote}
(i) 中原文作「对 $X$ 的每个子集 $U$」，疑为「每个开子集」之误：前层只在开子集上有定义，且紧随其后的 (ii) 明确写「开子集」。译文照录原文，读者按「开子集」理解即可。[依据：原书 p.12 推论 1.1.8 的 (i) 与 (ii)]
\end{tnote}

我们留给读者证明以下推论。

\begin{corollary}\label{it:1.1.9}
设 $X$ 是拓扑空间。
\begin{enumerate}[label=(\roman*),leftmargin=2.8em]
  \item 在 $X$ 上交换群的前层范畴中，序列
  \[
  \cF\xrightarrow{\phi}\cG\xrightarrow{\psi}\cH
  \]
  正合，当且仅当
  \[
  \cF(U)\xrightarrow{\phi(U)}\cG(U)\xrightarrow{\psi(U)}\cH(U)
  \]
  对 $X$ 的每个开子集 $U$ 正合。
  \item 在 $X$ 上交换群的层范畴中，序列
  \[
  \cF\xrightarrow{\phi}\cG\xrightarrow{\psi}\cH
  \]
  正合，当且仅当
  \[
  \cF_P\xrightarrow{\phi_P}\cG_P\xrightarrow{\psi_P}\cH_P
  \]
  对任意 $P\in X$ 正合。
\end{enumerate}
\end{corollary}
\begin{tnote}
\ref{it:1.1.9} 的 (ii) 与 (i) 的对比是本节要害：前层的正合可以逐个开集检验，层的正合却不能——层的满射只是「局部提升截面」，并不保证 $\cG(U)\to\cH(U)$ 满（见 \ref{it:1.1.8}(iv)），所以层的正合性只能用茎来刻画。[依据：原书 p.12–13 推论 1.1.9 与 p.12 推论 1.1.8 (iv)]
\end{tnote}

设 $f\colon X\to Y$ 是拓扑空间的连续映射。对 $X$ 上的每个层 $\cF$，定义 $\cF$ 的\emph{直接像}\idx{zhijiexiang}{直接像} $f_*\cF$ 为 $Y$ 上由
\[
(f_*\cF)(V)=\cF(f^{-1}(V))
\]
定义的层，其中 $V$ 取遍 $Y$ 的开子集。注意 $f_*$ 是从 $X$ 上交换群的层范畴到 $Y$ 上交换群的层范畴的加性函子。对 $Y$ 上的每个层 $\cG$，定义 $\cG$ 的\emph{原像}\idx{yuanxiang}{原像} $f^{-1}\cG$ 为 $X$ 上与前层
\[
U\mapsto\dlim_{f(U)\subset V}\cG(V)
\]
所伴随的层，其中 $U$ 取遍 $X$ 的开子集，而有向极限取遍 $Y$ 中包含 $f(U)$ 的开子集 $V$ 构成的族。注意 $f^{-1}$ 是从 $Y$ 上交换群的层范畴到 $X$ 上交换群的层范畴的加性函子。当 $f$ 是嵌入时，我们常把 $f^{-1}\cG$ 记为 $\cG|_X$，并称它为 $\cG$ 在 $X$ 上的\emph{限制}。

对 $X$ 上的每个层 $\cF$ 与 $Y$ 上的每个层 $\cG$，有一一对应
\[
\Hom(\cG,f_*\cF)\cong\Hom(f^{-1}\cG,\cF).
\]
事实上，可以验证，集合 $\Hom(\cG,f_*\cF)$ 与集合 $\Hom(f^{-1}\cG,\cF)$ 都一一对应于由元组 $(\phi_{U,V})_{(U,V)}$ 构成的集合，其中 $(U,V)$ 取遍所有满足以下条件的对：$U$ 是 $X$ 的开子集，$V$ 是 $Y$ 的开子集，且 $f(U)\subset V$；而 $\phi_{U,V}\colon\cG(V)\to\cF(U)$ 是同态，使得若 $(U',V')$ 是另一个这样的对且 $U'\subset U$、$V'\subset V$，则下图交换：
\begin{center}
\begin{tikzcd}
\cG(V) \arrow[r,"\phi_{U,V}"] \arrow[d,"\rho_{VV'}"'] & \cF(U) \arrow[d,"\rho_{UU'}"] \\
\cG(V') \arrow[r,"\phi_{U'V'}"'] & \cF(U').
\end{tikzcd}
\end{center}

一一对应 $\Hom(\cG,f_*\cF)\cong\Hom(f^{-1}\cG,\cF)$ 对 $\cF$ 与 $\cG$ 具有函子性，也就是说，对 $X$ 上层之间的每个态射 $\cF_1\to\cF_2$ 与 $Y$ 上层之间的每个态射 $\cG_1\to\cG_2$，下图交换：
\begin{center}
\begin{tikzcd}
\Hom(\cG_2,f_*\cF_1) \arrow[r,"\cong"] \arrow[d] & \Hom(f^{-1}\cG_2,\cF_1) \arrow[d] \\
\Hom(\cG_1,f_*\cF_2) \arrow[r,"\cong"'] & \Hom(f^{-1}\cG_1,\cF_2),
\end{tikzcd}
\end{center}
其中竖直箭头由 $\cF_1\to\cF_2$ 与 $\cG_1\to\cG_2$ 诱导。

一般地，设 $\cC$ 与 $\cD$ 是范畴，$u\colon\cC\to\cD$ 与 $v\colon\cD\to\cC$ 是函子。若对 $\cC$ 中每个对象 $C$ 与 $\cD$ 中每个对象 $D$，都有双射
\[
\alpha_{C,D}\colon\Hom(C,v(D))\cong\Hom(u(C),D),
\]
且它对 $C$ 与 $D$ 具有函子性，则称 $u$ \emph{左伴随}于 $v$，或称 $v$ \emph{右伴随}于 $u$。我们刚刚看到，$\cG\mapsto f^{-1}\cG$ 左伴随于 $\cF\mapsto f_*\cF$。\ref{it:1.1.5} 说明，从层的范畴到前层的范畴的包含函子的左伴随是从前层范畴到层范畴的函子 $\cF\mapsto\cF^+$。

若左伴随于 $v$ 的函子存在，则它在同构意义下唯一。设 $u$ 与 $u'$ 是左伴随于 $v$ 的两个函子。有双射
\[
\Hom(u(C),D)\cong\Hom(u'(C),D),
\]
它对 $C$ 与 $D$ 具有函子性。这些双射由对 $C$ 具有函子性的同构 $u'(C)\cong u(C)$ 诱导。事实上，取 $D=u(C)$，设 $\phi_C\in\Hom(u'(C),u(C))$ 是对应于 $\id_{u(C)}\in\Hom(u(C),u(C))$ 的态射。取 $D=u'(C)$，设 $\psi_C\in\Hom(u(C),u'(C))$ 是对应于 $\id_{u'(C)}\in\Hom(u'(C),u'(C))$ 的态射。可以验证，$\phi$ 与 $\psi$ 是 $u$ 与 $u'$ 之间的自然变换，且互为逆。

类似地，若右伴随于 $u$ 的函子存在，则它在同构意义下唯一。

仍然假设有双射
\[
\alpha_{C,D}\colon\Hom(C,v(D))\cong\Hom(u(C),D),
\]
且它对 $C$ 与 $D$ 具有函子性。取 $C=v(D)$，设 $\mathrm{adj}_D\colon uv(D)\to D$ 是 $\Hom(uv(D),D)$ 中对应于 $\id_{v(D)}$ 的态射。则 $\mathrm{adj}$ 是从函子 $uv$ 到函子 $\id_{\cD}$ 的自然变换，称为\emph{伴随态射}\idx{bansuitaishe}{伴随态射}。对 $\Hom(C,v(D))$ 中任意元素 $\phi\colon C\to v(D)$，$\alpha_{C,D}(\phi)$ 等于复合
\[
u(C)\xrightarrow{u(\phi)}uv(D)\xrightarrow{\mathrm{adj}_D}D.
\]
这可由以下交换图得到：
\begin{center}
\begin{tikzcd}
\Hom(C,v(D)) \arrow[r,"\alpha_{C,D}"] & \Hom(u(C),D) \\
\Hom(v(D),v(D)) \arrow[r,"\alpha_{v(D),D}"'] \arrow[u] & \Hom(uv(D),D) \arrow[u]
\end{tikzcd}
\end{center}
其中竖直箭头由与 $\phi\colon C\to v(D)$ 及 $u(\phi)\colon u(C)\to uv(D)$ 的复合诱导。

取 $D=u(C)$，设 $\mathrm{adj}_C\colon C\to vu(C)$ 是 $\Hom(C,vu(C))$ 中对应于 $\id_{u(C)}$ 的态射。则 $\mathrm{adj}$ 是从函子 $\id_{\cC}$ 到函子 $vu$ 的自然变换，也称为\emph{伴随态射}。对 $\Hom(u(C),D)$ 中任意元素 $\psi\colon u(C)\to D$，$\alpha^{-1}_{C,D}(\psi)$ 等于复合
\[
C\xrightarrow{\mathrm{adj}_C}vu(C)\xrightarrow{v(\psi)}v(D).
\]

设 $f\colon X\to Y$ 是拓扑空间的连续映射。对 $X$ 上的任意层 $\cF$，如下定义一个典范态射 $f^{-1}f_*\cF\to\cF$。因 $f^{-1}f_*\cF$ 是前层
\[
U\mapsto\dlim_{f(U)\subset V}(f_*\cF)(V)=\dlim_{U\subset f^{-1}(V)}\cF(f^{-1}(V))
\]
所伴随的层，由 \ref{it:1.1.5}，只需定义从前层 $U\mapsto\dlim_{U\subset f^{-1}(V)}\cF(f^{-1}(V))$ 到层 $\cF$ 的一个典范态射。我们定义
\[
\dlim_{U\subset f^{-1}(V)}\cF(f^{-1}(V))\to\cF(U)
\]
为限制 $\cF(f^{-1}(V))\to\cF(U)$ 诱导的映射。

对 $Y$ 上的任意层 $\cG$，如下定义一个典范态射 $\cG\to f_*f^{-1}\cG$。对 $Y$ 的任意开子集 $W$，因 $f(f^{-1}(W))\subset W$，有映射
\[
\cG(W)\to\dlim_{f(f^{-1}(W))\subset V}\cG(V).
\]
把它与典范同态
\[
\dlim_{f(f^{-1}(W))\subset V}\cG(V)\to f^{-1}\cG(f^{-1}(W))
\]
复合，得到同态 $\cG(W)\to f^{-1}\cG(f^{-1}(W))=(f_*f^{-1}\cG)(W)$，从而得到层态射 $\cG\to f_*f^{-1}\cG$。

典范态射 $\cG\to f_*f^{-1}\cG$ 与 $f^{-1}f_*\cF\to\cF$ 是伴随函子对 $(f^{-1},f_*)$ 的伴随态射。定义
\[
\alpha_{\cF,\cG}\colon\Hom(\cG,f_*\cF)\to\Hom(f^{-1}\cG,\cF)
\]
如下：对每个态射 $\phi\colon\cG\to f_*\cF$，$\alpha_{\cF,\cG}(\phi)$ 是复合
\[
f^{-1}\cG\xrightarrow{f^{-1}\phi}f^{-1}f_*\cF\to\cF.
\]
定义
\[
\beta_{\cF,\cG}\colon\Hom(f^{-1}\cG,\cF)\to\Hom(\cG,f_*\cF)
\]
如下：对每个态射 $\psi\colon f^{-1}\cG\to\cF$，$\beta_{\cF,\cG}(\psi)$ 是复合
\[
\cG\to f_*f^{-1}\cG\xrightarrow{f_*\psi}f_*\cF.
\]
则 $\alpha_{\cF,\cG}$ 与 $\beta_{\cF,\cG}$ 互为逆。这给出了一一对应 $\Hom(\cG,f_*\cF)\cong\Hom(f^{-1}\cG,\cF)$ 的另一种描述。

设 $f\colon X\to Y$ 与 $g\colon Y\to Z$ 是两个连续映射。显然有 $(gf)_*=g_*f_*$。于是对 $X$ 上的任意层 $\cF$ 与 $Z$ 上的任意层 $\cH$，有
\begin{align*}
\Hom((gf)^{-1}\cH,\cF)&\cong\Hom(\cH,(gf)_*\cF)\cong\Hom(\cH,g_*f_*\cF)\\
&\cong\Hom(g^{-1}\cH,f_*\cF)\cong\Hom(f^{-1}g^{-1}\cH,\cF),
\end{align*}
即
\[
\Hom((gf)^{-1}\cH,\cF)\cong\Hom(f^{-1}g^{-1}\cH,\cF).
\]
这说明 $(gf)^{-1}\cong f^{-1}g^{-1}$。

设 $P$ 是 $X$ 中的点，$i\colon\{P\}\to X$ 是包含映射。则 $X$ 上层 $\cF$ 的原像 $i^{-1}\cF$ 可与茎 $\cF_P$ 等同。对每个连续映射 $f\colon X\to Y$ 与 $Y$ 上的每个层 $\cG$，有
\[
i^{-1}f^{-1}\cG\cong(fi)^{-1}\cG.
\]
于是 $(f^{-1}\cG)_P\cong\cG_{f(P)}$。

本节以两个在任何阿贝尔范畴中都成立的引理结束。我们只需要阿贝尔范畴是层范畴这一特殊情形；此时这两个引理可以容易地用图追踪证明。

\begin{lemma}[（蛇引理\idx{sheyinli}{蛇引理}）]\label{it:1.1.10}
设
\begin{center}
\begin{tikzcd}
0 \arrow[r] & A \arrow[r] \arrow[d,"u"] & B \arrow[r] \arrow[d,"v"] & C \arrow[r] \arrow[d,"w"] & 0 \\
0 \arrow[r] & A' \arrow[r] & B' \arrow[r] & C' \arrow[r] & 0
\end{tikzcd}
\end{center}
是阿贝尔范畴中的交换图，且两行都正合。则存在态射 $\delta\colon\ker w\to\coker u$，使得以下序列正合：
\[
0\to\ker u\to\ker v\to\ker w\xrightarrow{\delta}\coker u\to\coker v\to\coker w\to 0.
\]
\end{lemma}

\begin{lemma}[（五引理\idx{wuyinli}{五引理}）]\label{it:1.1.11}
设
\begin{center}
\begin{tikzcd}
A_1 \arrow[r] \arrow[d,"f_1"] & A_2 \arrow[r] \arrow[d,"f_2"] & A_3 \arrow[r] \arrow[d,"f_3"] & A_4 \arrow[r] \arrow[d,"f_4"] & A_5 \arrow[d,"f_5"] \\
B_1 \arrow[r] & B_2 \arrow[r] & B_3 \arrow[r] & B_4 \arrow[r] & B_5
\end{tikzcd}
\end{center}
是阿贝尔范畴中的交换图，且两行都正合。若 $f_1$、$f_2$、$f_4$ 与 $f_5$ 都是同构，则 $f_3$ 也是同构。
\end{lemma}

\begin{exercises}
\item \label{ex:1.1.1} \emph{过滤交换群}\idx{guolvjiaohuanqun}{过滤交换群}是一个对 $(A,F^{\bullet}A)$，其中 $A$ 是交换群，且
\[
\cdots\supset F^iA\supset F^{i+1}A\supset\cdots
\]
是指标 $i\in\bZ$ 的一族递减子群。过滤交换群之间的同态 $f\colon(A,F^{\bullet}A)\to(B,F^{\bullet}B)$ 是交换群同态 $f\colon A\to B$，使 $f(F^iA)\subset F^iB$ 对所有 $i$ 成立。证明：过滤交换群构成有零对象、且每个态射都有核与余核的加性范畴。给出一个态射 $f$ 的例子，使典范态射 $\operatorname{coim}f\to\operatorname{im}f$ 不是同构。

\item \label{ex:1.1.2} 设 $X$ 是拓扑空间，$A$ 是赋以离散拓扑的交换群，$\cF$ 是这样的层：对 $X$ 的每个开子集 $U$，$\cF(U)$ 是从 $U$ 到 $A$ 的连续映射构成的群。证明：$\cF$ 同构于与常前层 $U\mapsto A$ 相伴的层。

\item \label{ex:1.1.3} 对复平面 $\bC$ 的每个开子集 $U$，设 $\cO(U)$ 是 $U$ 上全纯函数构成的环，$\cO^*(U)$ 是 $\cO(U)$ 中的单位群。证明：由
\[
\cO(U)\to\cO^*(U),\quad f\mapsto e^{2\pi if}
\]
定义的态射 $\cO\to\cO^*$ 在交换群的层范畴中是满态射，但在前层范畴中不是满态射。这里把 $\cO$ 视为关于函数加法的交换群层。证明：该态射的核同构于与常前层 $U\mapsto\bZ$ 相伴的层。

\item \label{ex:1.1.4} 设 $0\to A\to B\to C\to 0$ 是阿贝尔范畴中的态射序列。证明以下陈述等价：
\begin{enumerate}[label=(\arabic*),leftmargin=2.4em]
  \item 该序列是短正合序列。
  \item $B\to C$ 是满态射，且 $A\to B$ 是它的核。
  \item $A\to B$ 是单态射，且 $B\to C$ 是它的余核。
\end{enumerate}

\item \label{ex:1.1.5} 设 $A$ 与 $B$ 是阿贝尔范畴中的对象。证明典范序列
\[
0\to A\to A\oplus B\to B\to 0
\]
正合。

\item \label{ex:1.1.6} 设 $I$ 是对象构成集合的范畴，$F$ 是从 $I$ 到交换群范畴的共变函子。对每个 $i\in I$，设 $k_i\colon F(i)\to\bigoplus_{i\in I}F(i)$ 是典范单态射。设 $H$ 是 $\bigoplus_{i\in I}F(i)$ 中由
\[
k_i(x_i)-k_j(F(i\to j)(x_i)),
\]
生成的子群，其中 $i\to j$ 取遍 $I$ 中所有态射，$x_i$ 取遍 $F(i)$ 的所有元素。令
\[
\dlim_{i\in I}F(i)=\Bigl(\bigoplus_{i\in I}F(i)\Bigr)/H.
\]
设 $\phi_i\colon F(i)\to\dlim_{i\in I}F(i)$ 是 $k_i$ 与投影 $\bigoplus_{i\in I}F(i)\to\bigl(\bigoplus_{i\in I}F(i)\bigr)/H$ 的复合。对 $I$ 中每个态射 $i\to j$，有 $\phi_j\circ F(i\to j)=\phi_i$。若 $A$ 是交换群，$\psi_i\colon F(i)\to A$（$i\in I$）是一族同态，使 $\psi_j\circ F(i\to j)=\psi_i$ 对 $I$ 中所有态射 $i\to j$ 成立，则存在唯一一个同态 $\psi\colon\dlim_{i\in I}F(i)\to A$，使 $\psi\phi_i=\psi_i$ 对所有 $i$ 成立。

\item \label{ex:1.1.7} 设 $(A_i,\phi_{ji})_{i\in\bN}$ 是自然数的有向集 $(\bN,\le)$ 上的交换群逆向系统。考虑同态
\[
f\colon\prod_{i\in\bN}A_i\to\prod_{i\in\bN}A_i,\quad (a_i)\mapsto(a_i-\phi_{i+1,i}(a_{i+1})).
\]
定义
\[
\plim_i^1A_i=\coker f.
\]
设 $u\colon(A'_i,\phi'_{ji})_{i\in\bN}\to(A_i,\phi_{ji})_{i\in\bN}$ 与 $v\colon(A_i,\phi_{ji})_{i\in\bN}\to(A''_i,\phi''_{ji})_{i\in\bN}$ 是交换群逆向系统的态射，使得序列
\[
0\to A'_i\xrightarrow{u_i}A_i\xrightarrow{v_i}A''_i\to 0
\]
对所有 $i$ 正合。证明：有正合序列
\[
0\to\plim_i A'_i\to\plim_i A_i\to\plim_i A''_i\to\plim_i^1A'_i\to\plim_i^1A_i\to\plim_i^1A''_i\to 0.
\]

\item \label{ex:1.1.8} 设 $\cC$ 是范畴。对任意对象 $X\in\ob\cC$，设 $\widetilde X\colon\cC\to(\mathrm{Sets})$ 是由
\[
\widetilde X(Y)=\Hom(Y,X)
\]
定义的反变函子。若一个从 $\cC$ 到集合范畴的函子同构于 $\widetilde X$，则称它可由 $X$ 表出。证明：对任意反变函子 $G\colon\cC\to(\mathrm{Sets})$，有一一对应
\[
\Hom(\widetilde X,G)\to G(X),\quad \alpha\mapsto\alpha_X(\id_X),
\]
其中 $\Hom(\widetilde X,G)$ 是从函子 $\widetilde X$ 到函子 $G$ 的自然变换构成的集合。对共变函子证明同样的结果。

\item \label{ex:1.1.9} 设 $u\colon\cC\to\cD$ 是函子。假设对每个对象 $D\in\ob\cC$，函子
\[
\cC\to(\mathrm{Sets}),\quad C\mapsto\Hom(u(C),D)
\]
可由某个对象 $v(D)\in\ob\cC$ 表出。则 $v\colon\cD\to\cC$ 是 $u$ 的右伴随函子。

\item \label{ex:1.1.10} 设 $u\colon\cC\to\cD$ 是函子。若对任意对象 $C_1,C_2\in\ob\cC$，映射
\[
\Hom(C_1,C_2)\to\Hom(u(C_1),u(C_2))
\]
是单射（resp.\ 双射），则称 $u$ 是\emph{忠实}（resp.\ \emph{全忠实}）的。若对 $\cD$ 中任意对象 $D$，存在 $\cC$ 中的对象 $C$，使得有同构 $u(C)\cong D$，则称 $u$ 是\emph{本质满}的。若 $u$ 既全忠实又本质满，则称 $u$ 是\emph{范畴等价}。假设 $u$ 是范畴等价。对任意对象 $D\in\ob\cC$，选取对象 $v(D)\in\ob\cC$ 使 $uv(D)\cong D$。证明：$v$ 是函子，并且它同时是 $u\colon\cD\to\cC$ 的左伴随与右伴随。它称为 $u$ 的\emph{准逆}。

\begin{tnote}
习题 10 原书印作「对任意对象 $D\in\ob\cC$，选取对象 $v(D)\in\ob\cC$ 使 $uv(D)\cong D$」；因 $u\colon\cC\to\cD$，此处 $D$ 应在 $\cD$ 中（$uv(D)$ 才与 $D$ 同属 $\cD$），原文疑为笔误（应作 $D\in\ob\cD$）。译文照原书印法照录。[依据：原书 p.19 习题 10]
\end{tnote}
\end{exercises}
